% ============================================================
%  PRUEBA TÉCNICA — SISTEMAS DE ACOPLAMIENTO ELECTROMAGNÉTICO
%  Formato: Technical Report / IEEE-style
% ============================================================
\documentclass[11pt,a4paper]{article}

% ─── Codificación y fuentes ───────────────────────────────
\usepackage[T1]{fontenc}
\usepackage[utf8]{inputenc}
\usepackage[spanish,activeacute]{babel}
\usepackage{mathpazo}   % Palatino + Euler math
\usepackage{microtype}

% ─── Matemáticas ──────────────────────────────────────────
\usepackage{amsmath,amssymb,amsthm}
\usepackage{mathtools}
\usepackage{siunitx}
\sisetup{output-decimal-marker={,}, separate-uncertainty=true}

% ─── Página y márgenes ────────────────────────────────────
\usepackage[
  a4paper,
  top=25mm, bottom=28mm,
  left=22mm, right=22mm,
  headheight=14pt
]{geometry}

% ─── Color ────────────────────────────────────────────────
\usepackage[dvipsnames,table]{xcolor}
\definecolor{azulprincipal}{RGB}{0,74,140}
\definecolor{azulclaro}{RGB}{220,235,250}
\definecolor{azulmed}{RGB}{180,210,240}
\definecolor{gristexto}{RGB}{60,60,70}
\definecolor{grisinfobox}{RGB}{245,247,250}
\definecolor{naranja}{RGB}{190,80,20}
\definecolor{verdemed}{RGB}{30,115,60}
\definecolor{rojoaviso}{RGB}{160,30,20}
\definecolor{grisbordes}{RGB}{180,185,195}

% ─── Encabezados y pie de página ──────────────────────────
\usepackage{fancyhdr}
\pagestyle{fancy}
\fancyhf{}
\fancyhead[L]{\small\color{gristexto}\textsc{Prueba Técnica de Ingeniería}}
\fancyhead[R]{\small\color{gristexto}\textsc{Acoplamiento Electromagnético}}
\fancyfoot[C]{\small\color{gristexto}\thepage}
\renewcommand{\headrulewidth}{0.4pt}
\renewcommand{\headrule}{\hbox to\headwidth{\color{azulprincipal}\leaders\hrule height \headrulewidth\hfill}}
\renewcommand{\footrulewidth}{0pt}

% ─── Títulos de sección ───────────────────────────────────
\usepackage{titlesec}
\titleformat{\section}
  {\large\bfseries\color{azulprincipal}}
  {\thetitle.\quad}{0pt}{}[{\color{azulprincipal}\titlerule[0.6pt]}]
\titlespacing{\section}{0pt}{18pt}{8pt}

\titleformat{\subsection}
  {\normalsize\bfseries\color{gristexto}}
  {\thetitle\quad}{0pt}{}
\titlespacing{\subsection}{0pt}{12pt}{5pt}

\titleformat{\subsubsection}
  {\small\bfseries\itshape\color{naranja}}
  {\thetitle\quad}{0pt}{}
\titlespacing{\subsubsection}{0pt}{8pt}{3pt}

% ─── Tablas y figuras ─────────────────────────────────────
\usepackage{booktabs}
\usepackage{array}
\usepackage{multirow}
\usepackage{tabularx}
\usepackage{longtable}
\usepackage{graphicx}
\usepackage{caption}
\captionsetup{
  font={small,it},
  labelfont={bf,color=azulprincipal},
  labelsep=period,
  skip=6pt
}

% ─── Cajas y entornos especiales ──────────────────────────
\usepackage{tcolorbox}
\tcbuselibrary{skins,breakable,theorems}

% Caja de contexto / descripción
\newtcolorbox{contextobox}[1][]{
  enhanced, breakable,
  colback=grisinfobox,
  colframe=azulprincipal,
  boxrule=0.5pt,
  leftrule=3pt,
  arc=2pt,
  fonttitle=\small\bfseries\color{azulprincipal},
  title={Contexto},
  #1
}

% Caja de requisito
\newtcolorbox{requisitobox}[2][]{
  enhanced,
  colback=white,
  colframe=grisbordes,
  boxrule=0.5pt,
  leftrule=2.5pt,
  arc=1pt,
  left=6pt, right=6pt, top=4pt, bottom=4pt,
  fontupper=\small,
  overlay={
    \node[anchor=north west, fill=azulprincipal, text=white,
          font=\tiny\bfseries\ttfamily, inner sep=2pt, minimum width=1cm]
      at (frame.north west) {#2};
  },
  #1
}

% Caja de aviso ético
\newtcolorbox{avisobox}{
  enhanced, breakable,
  colback=yellow!6!white,
  colframe=rojoaviso,
  boxrule=0.5pt,
  leftrule=3pt,
  arc=2pt,
  fonttitle=\small\bfseries\color{rojoaviso},
  title={\faWarning~Nota Ética y Regulatoria},
}

% Caja de fórmula
\newtcolorbox{formulabox}[1][]{
  enhanced,
  colback=azulclaro!40!white,
  colframe=azulprincipal!60!white,
  boxrule=0.4pt,
  arc=2pt,
  left=10pt, right=10pt, top=6pt, bottom=6pt,
  #1
}

% Caja de tarea
\newtcolorbox{tareabox}[2][]{
  enhanced, breakable,
  colback=white,
  colframe=azulprincipal,
  boxrule=0.5pt,
  arc=2pt,
  fonttitle=\small\bfseries,
  coltitle=white,
  attach boxed title to top left={yshift=-\tcboxedtitleheight/2, xshift=6pt},
  boxed title style={colback=azulprincipal, arc=1pt, boxrule=0pt},
  title={#2},
  #1
}

% Caja de criterio
\newtcolorbox{criteriobox}[1][]{
  enhanced,
  colback=verdemed!5!white,
  colframe=verdemed!60!white,
  boxrule=0.4pt,
  leftrule=2.5pt,
  arc=1pt,
  fonttitle=\small\bfseries\color{verdemed},
  #1
}

% ─── Listas personalizadas ────────────────────────────────
\usepackage{enumitem}
\setlist[itemize]{
  leftmargin=1.2em,
  itemsep=2pt,
  topsep=4pt,
  label={\color{azulprincipal}\small$\blacktriangleright$}
}
\setlist[enumerate]{
  leftmargin=1.4em,
  itemsep=3pt,
  topsep=4pt,
  label={\color{azulprincipal}\textbf{\arabic*.}}
}

% ─── TikZ para diagramas ──────────────────────────────────
\usepackage{tikz}
\usetikzlibrary{shapes.geometric, arrows.meta, positioning, decorations.pathmorphing, backgrounds, fit, calc, shadows}

\tikzset{
  bloque/.style={
    rectangle, rounded corners=2pt, minimum width=2.4cm, minimum height=0.9cm,
    text centered, font=\small\bfseries, draw=azulprincipal, fill=azulclaro,
    text=gristexto, general shadow={shadow scale=1.005, shadow xshift=0.5pt, shadow yshift=-0.5pt, fill=black!20, opacity=0.5}
  },
  bloquenaranja/.style={bloque, draw=naranja, fill=naranja!10},
  bloqueverde/.style={bloque, draw=verdemed, fill=verdemed!10},
  flecha/.style={-Stealth, thick, color=azulprincipal},
  flechadouble/.style={Stealth-Stealth, thick, color=naranja, dashed},
}

% ─── Hipervínculos ────────────────────────────────────────
\usepackage{hyperref}
\hypersetup{
  colorlinks=true,
  linkcolor=azulprincipal,
  urlcolor=azulprincipal,
  citecolor=naranja,
  pdftitle={Prueba Técnica -- Acoplamiento Electromagnético},
  pdfauthor={Ingeniería Electromagnética},
  pdfsubject={Technical Test},
}

% ─── Utilidades ───────────────────────────────────────────
\usepackage{parskip}
\setlength{\parskip}{5pt}
\usepackage{lipsum}

% Operadores matemáticos
\DeclareMathOperator{\sinc}{sinc}
\newcommand{\vect}[1]{\boldsymbol{#1}}

% ─── Línea divisoria de sección ───────────────────────────
\newcommand{\partedivider}[1]{%
  \vspace{10pt}
  \noindent
  \begin{tikzpicture}[remember picture]
    \fill[azulprincipal] (0,0) rectangle (\linewidth, 3pt);
    \node[anchor=west, fill=azulprincipal, text=white,
          font=\large\bfseries, inner xsep=8pt, inner ysep=6pt]
      at (0,-0.05) {#1};
  \end{tikzpicture}
  \vspace{6pt}
}

% ============================================================
\begin{document}
% ============================================================

% ─── PORTADA ─────────────────────────────────────────────
\thispagestyle{empty}
\begin{titlepage}

\begin{tikzpicture}[remember picture, overlay]
  % Banda superior
  \fill[azulprincipal] (current page.north west) rectangle
    ([yshift=-4cm] current page.north east);
  % Detalle geométrico
  \fill[azulprincipal!60!black] (current page.north east) --
    ([xshift=-6cm, yshift=-4cm] current page.north east) --
    ([yshift=-4cm] current page.north east) -- cycle;
  % Banda inferior
  \fill[azulprincipal!20!white] (current page.south west) rectangle
    ([yshift=1.4cm] current page.south east);
  \node[anchor=south east, font=\small\color{azulprincipal}, inner sep=8pt]
    at ([yshift=1.4cm] current page.south east)
    {\textit{Versión para evaluación académica}};
\end{tikzpicture}

\vspace*{1.8cm}

{\raggedright
  {\color{white}\fontsize{9}{11}\selectfont\bfseries\MakeUppercase{Prueba Técnica de Ingeniería}}
  \par\vspace{0.5cm}
  {\color{white}\fontsize{24}{28}\selectfont\bfseries
    Sistemas de Acoplamiento\\[4pt]
    Electromagnético
  }
  \par\vspace{0.4cm}
  {\color{azulmed}\fontsize{14}{18}\selectfont\itshape
    Diseño, Análisis y Validación Experimental
  }
}

\vspace{2.8cm}

\noindent
\begin{tabular}{@{}p{0.46\linewidth}p{0.46\linewidth}@{}}
  \toprule[1.2pt]
  \multicolumn{2}{l}{\small\bfseries\color{azulprincipal} METADATOS DEL DOCUMENTO} \\[4pt]
  \small\textbf{Disciplina} & \small Ingeniería Electromagnética \\[2pt]
  \small\textbf{Modalidad}  & \small Teórico-Práctica \\[2pt]
  \small\textbf{Nivel}      & \small Ingeniería Avanzada \\[2pt]
  \small\textbf{Extensión}  & \small Máximo 15 páginas \\[2pt]
  \small\textbf{Formato}    & \small Documento técnico + demostración \\[2pt]
  \bottomrule[1.2pt]
\end{tabular}

\vspace{1.4cm}

\noindent\textbf{\color{azulprincipal}Resumen ejecutivo.}
Esta prueba técnica evalúa simultáneamente el dominio teórico de los principios
electromagnéticos y las habilidades prácticas de diseño e implementación de circuitos.
Consta de dos partes diferenciadas: una sección teórica orientada al análisis de sistemas
de pulsos electromagnéticos de alta frecuencia, y una sección práctica centrada en la
demostración experimental del acoplamiento inductivo resonante entre circuitos de baja
tensión. Se valoran especialmente el rigor analítico, la calidad de la justificación técnica
y la capacidad de relacionar fenómenos físicos con limitaciones reales de hardware.

\vfill

\noindent\begin{minipage}{\linewidth}
  \small\color{gray}
  \textbf{Nota:} Este documento tiene carácter estrictamente académico y formativo.
  Ninguna sección del mismo constituye una guía operativa para el desarrollo de
  dispositivos con fines no autorizados.
\end{minipage}

\end{titlepage}

% ─── TABLA DE CONTENIDOS ────────────────────────────────
\setcounter{page}{1}
\tableofcontents
\vspace{10pt}
\noindent\color{grisbordes}\rule{\linewidth}{0.4pt}
\newpage

% ============================================================
\partedivider{PARTE I \quad Análisis Teórico: Pulsos Electromagnéticos}
% ============================================================

\section{Contexto y Especificaciones del Sistema}

\begin{contextobox}
Se solicita el diseño conceptual de un sistema de pulsos electromagnéticos (EMP, del
inglés \textit{Electromagnetic Pulse}) orientado a la inhabilitación de sistemas electrónicos
embarcados. El sistema debe emitir pulsos de alta energía en un rango de frecuencias que
comienza en $f_{\min} = \SI{300}{\mega\hertz}$, buscando maximizar el espectro de cobertura
por encima de dicha frecuencia de corte. El diseño contempla acoplamiento efectivo a
distancia variable~$d$ y adaptabilidad a diferentes frecuencias de resonancia del objetivo.
\end{contextobox}


\subsection{Requisitos del sistema}

\begin{requisitobox}{R-01}
\textbf{Rango de frecuencia.} El sistema debe emitir pulsos efectivos desde
$f_{\min} = \SI{300}{\mega\hertz}$ hasta la mayor cobertura espectral posible.
La relación entre la anchura temporal del pulso~$\tau_p$ y su contenido espectral viene dada por:
\[
  \Delta f \approx \frac{1}{\tau_p}
\]
lo que implica que pulsos más cortos generan mayor cobertura en frecuencia.
\end{requisitobox}

\vspace{6pt}

\begin{requisitobox}{R-02}
\textbf{Acoplamiento a distancia.} El sistema debe garantizar acoplamiento efectivo
con objetivos a distancia variable~$d$. Debe diferenciarse entre la región de campo
cercano ($d < \lambda/2\pi$) y la de campo lejano ($d \gg \lambda$), dado que las
leyes de atenuación difieren sustancialmente en cada zona.
\end{requisitobox}

\vspace{6pt}

\begin{requisitobox}{R-03}
\textbf{Maximización de energía transmitida.} Optimizar la energía entregada al
objetivo, considerando las limitaciones térmicas, la ganancia de antena disponible
y los márgenes de susceptibilidad del objetivo.
\end{requisitobox}

\vspace{6pt}

\begin{requisitobox}{R-04}
\textbf{Adaptabilidad espectral.} El sistema debe operar eficazmente con circuitos
objetivo cuya frecuencia de resonancia~$f_0$ no está predeterminada, ya sea mediante
técnicas de banda ancha (UWB) o barrido dinámico de frecuencia.
\end{requisitobox}

% ============================================================
\section{Tareas a Realizar --- Parte Teórica}

\subsection{T-01: Diseño conceptual del sistema}

Se propone el siguiente \textbf{diagrama de bloques} del sistema completo:

\vspace{8pt}
\begin{figure}[h!]
\centering
\begin{tikzpicture}[node distance=0.7cm and 0.5cm]

  % Nodos
  \node[bloque, fill=naranja!15, draw=naranja] (fuente) {Fuente de\\alimentación};
  \node[bloque, right=of fuente] (genp) {Generador\\de pulsos};
  \node[bloque, right=of genp]   (etapa) {Etapa de\\potencia};
  \node[bloque, right=of etapa]  (antena) {Sistema\\de antena};
  \node[bloque, fill=rojoaviso!10, draw=rojoaviso, right=of antena, xshift=0.4cm]
    (objetivo) {Objetivo\\(electrónica)};

  % Control (debajo)
  \node[bloquenaranja, below=1.2cm of genp, xshift=1.5cm]
    (control) {Unidad de\\control/sync.};
  \node[bloqueverde, below=1.2cm of etapa, xshift=1.4cm]
    (blindaje) {Blindaje\\EMP propio};

  % Flechas principales
  \draw[flecha] (fuente) -- (genp);
  \draw[flecha] (genp)   -- (etapa);
  \draw[flecha] (etapa)  -- (antena);
  \draw[flecha] (antena) -- node[above, font=\tiny, color=gray]{radiación} (objetivo);

  % Flechas control
  \draw[-Stealth, dashed, color=naranja] (control) -- (genp);
  \draw[-Stealth, dashed, color=naranja] (control) -- (etapa);
  \draw[-Stealth, dashed, color=naranja] (control) -- (antena);

  % Retroalimentación blindaje
  \draw[-Stealth, dashed, color=verdemed] (blindaje) -- (etapa);
  \draw[-Stealth, dashed, color=verdemed] (blindaje) -| (antena);

  % Etiquetas de flujo
  \node[font=\tiny\it, color=gray, below=3pt of fuente] {Alta tensión};
  \node[font=\tiny\it, color=gray, below=3pt of genp]   {Pulsos ns};
  \node[font=\tiny\it, color=gray, below=3pt of etapa]  {Amplificación};
  \node[font=\tiny\it, color=gray, below=3pt of antena] {Radiación};

\end{tikzpicture}
\caption{Diagrama de bloques del sistema EMP. Las líneas discontinuas representan señales de control y gestión.}
\end{figure}

\subsection{T-02: Análisis matemático}

\subsubsection{Potencia radiada y ecuación de Friis modificada}

La potencia recibida efectiva por el objetivo a distancia~$d$ en espacio libre se expresa mediante
la \textbf{ecuación de Friis}:

\begin{formulabox}
\begin{equation}
  P_r = P_t \, G_t(\theta, \phi) \, G_r \left(\frac{\lambda}{4\pi d}\right)^2
  \label{eq:friis}
\end{equation}
\end{formulabox}

donde $P_t$ es la potencia transmitida, $G_t$ y $G_r$ son las ganancias de antena transmisora
y receptora (objetivo), $\lambda = c/f$ es la longitud de onda y $d$ la distancia de separación.

Para deshabilitar un circuito electrónico se requiere superar su umbral de susceptibilidad
$P_{\text{sus}}$ (típicamente en el rango de \SIrange{0.1}{100}{\milli\watt} según el componente):

\begin{equation}
  P_t \geq P_{\text{sus}} \cdot \frac{(4\pi d)^2}{G_t \, G_r \, \lambda^2}
  \label{eq:potencianec}
\end{equation}

\subsubsection{Atenuación con la distancia y régimen de campo}

La frontera entre campo cercano y campo lejano para una antena de apertura efectiva $D$ viene dada por:
\[
  d_{\text{Fraunhofer}} = \frac{2D^2}{\lambda}
\]

En el \textbf{campo lejano} la densidad de potencia decrece como $d^{-2}$ (ley del inverso
del cuadrado). En el \textbf{campo cercano inductivo} la dependencia es de orden $d^{-3}$,
con lo que la energía disponible se reduce mucho más rápidamente:

\begin{equation}
  S(d) \propto \begin{cases}
    d^{-2} & d \gg \lambda/2\pi \quad \text{(campo lejano)} \\
    d^{-3} & d \ll \lambda/2\pi \quad \text{(campo cercano reactivo)}
  \end{cases}
\end{equation}

\subsubsection{Contenido espectral del pulso y transformada de Fourier}

Para un pulso rectangular de amplitud $A$ y duración $\tau_p$ el espectro de energía es:
\begin{equation}
  |E(f)|^2 = A^2 \tau_p^2 \sinc^2(f \tau_p), \qquad \sinc(x) \triangleq \frac{\sin(\pi x)}{\pi x}
  \label{eq:espectro}
\end{equation}

El primer cero del espectro se produce en $f = 1/\tau_p$, por lo que reducir $\tau_p$ desplaza
la energía hacia frecuencias más altas. Para cubrir el espectro a partir de \SI{300}{\mega\hertz}:
\[
  \tau_p \lesssim \frac{1}{\SI{300}{\mega\hertz}} \approx \SI{3.3}{\nano\second}
\]

\subsection{T-03: Diseño de antena}

\begin{tareabox}{Criterios de selección de antena}
Para un sistema EMP de banda ancha ($f \geq \SI{300}{\mega\hertz}$), la antena debe satisfacer:
\begin{itemize}
  \item \textbf{Ancho de banda elevado:} cociente $f_{\max}/f_{\min} \gg 1$ (antena UWB o de apertura).
  \item \textbf{Alta directividad:} concentrar la energía en la dirección del objetivo.
  \item \textbf{Alta eficiencia de radiación:} minimizar las pérdidas óhmicas en la estructura.
  \item \textbf{Bajo rizado de la fase:} para preservar la integridad del pulso temporal.
\end{itemize}
\end{tareabox}

\vspace{6pt}

\noindent La Tabla~\ref{tab:antenas} presenta una comparativa de topologías candidatas:

\begin{table}[h!]
\centering
\caption{Comparativa de topologías de antena para sistemas de banda ancha}
\label{tab:antenas}
\small
\rowcolors{2}{azulclaro!30}{white}
\begin{tabularx}{\linewidth}{@{}l c c c X@{}}
  \toprule[1.2pt]
  \rowcolor{azulprincipal}
  \color{white}\textbf{Topología} &
  \color{white}\textbf{BW relativo} &
  \color{white}\textbf{Ganancia} &
  \color{white}\textbf{$\eta_\text{rad}$} &
  \color{white}\textbf{Observaciones} \\
  \midrule
  Bocina piramidal  & Moderado (1:3)   & \SIrange{10}{25}{\dBi}  & $>90\%$ & Buena directividad, voluminosa \\
  Vivaldi graduada  & Muy alto (1:20+) & \SIrange{5}{12}{\dBi}   & $>85\%$ & Excelente para UWB, compacta \\
  Helicoidal axial  & Moderado (1:2)   & \SIrange{8}{15}{\dBi}   & $>80\%$ & Polarización circular \\
  Bicónica          & Alto (1:10)      & $\sim\SI{3}{\dBi}$      & $>75\%$ & Omnidireccional, cobertura total \\
  \textbf{TEM horn} & \textbf{Muy alto (1:30+)} & \textbf{\SIrange{8}{18}{\dBi}} & \textbf{$>92\%$} & \textbf{Candidata óptima} \\
  \bottomrule[1.2pt]
\end{tabularx}
\end{table}

\noindent\textbf{Justificación:} La antena de cuerno TEM (\textit{Transverse Electromagnetic horn})
presenta el mejor compromiso entre ancho de banda ultrancho, directividad moderada-alta y
preservación de la fase del pulso, siendo ampliamente empleada en sistemas de diagnóstico EMC.

\subsection{T-04: Consideraciones prácticas}

\subsubsection{Gestión térmica}

La potencia disipada en la etapa de conmutación puede estimarse como:
\[
  P_{\text{dis}} = P_{\text{entrada}} \cdot (1 - \eta_{\text{RF}})
\]
Para disiparla se propone: refrigeración líquida de la etapa de potencia, interruptores basados en
SiC o GaN (alta temperatura de operación, baja resistencia en conducción), y monitorización
en tiempo real de temperatura con corte de seguridad.

\subsubsection{Blindaje del propio sistema (EMP hardening)}

Los componentes de control y la fuente de alimentación deben protegerse con:
\begin{itemize}
  \item Apantallamiento de la caja con lámina de cobre ($t \geq 3\delta_s$, donde $\delta_s$ es la profundidad de efecto piel a $f_{\min}$).
  \item Filtros de entrada LC de orden elevado en todas las líneas de alimentación.
  \item Fotoacopladores en señales de control para eliminar caminos de acoplamiento galvánico.
\end{itemize}

\subsubsection{Sintonización dinámica}

Para adaptarse a una frecuencia de resonancia desconocida del objetivo se propone un
\textbf{barrido de frecuencia lineal (chirp)}, cuya representación temporal es:
\begin{equation}
  s(t) = A(t)\cos\!\left(2\pi f_0 t + \pi \mu t^2\right), \quad
  \mu = \frac{f_{\max}-f_{\min}}{T_{\text{barrido}}}
  \label{eq:chirp}
\end{equation}
Alternativamente, una antena de banda ancha real (UWB) que cubra el espectro completo
de forma simultánea elimina la necesidad de sintonización.

% ============================================================
\section{Criterios de Evaluación --- Parte Teórica}

\begin{criteriobox}[title={Rúbrica de evaluación teórica}, fonttitle=\small\bfseries\color{verdemed}]
\begin{enumerate}
  \item \textbf{Rigor técnico del diseño:} consistencia física y matemática del sistema propuesto.
  \item \textbf{Justificación de decisiones:} cada elección de diseño debe argumentarse con criterios cuantitativos.
  \item \textbf{Originalidad de soluciones:} creatividad en el tratamiento de los desafíos planteados.
  \item \textbf{Comprensión de principios físicos:} evidencia de dominio de electromagnetismo y teoría de circuitos.
\end{enumerate}
\end{criteriobox}

\newpage
% ============================================================
\partedivider{PARTE II \quad Práctica: Acoplamiento Inductivo Resonante}
% ============================================================

\section{Contexto y Especificaciones Técnicas}

\begin{contextobox}[title={Descripción del experimento}]
Se desarrollará una aplicación de \textbf{baja tensión} ($V_{CC} \in [\SI{5}{\volt}, \SI{12}{\volt}]$)
que demuestre experimentalmente los principios del acoplamiento inductivo resonante entre dos
circuitos RLC. El objetivo es transferir energía de forma inalámbrica desde un circuito transmisor~(A)
hasta un circuito receptor~(B), evidenciando el encendido de un LED como indicador visual de
transferencia energética.
\end{contextobox}

\subsection{Especificaciones de los circuitos}

\subsubsection{Circuito transmisor A}

\begin{itemize}
  \item Conectado a una fuente de alimentación de baja tensión ($\SI{5}{\volt}$ -- $\SI{12}{\volt}$ DC).
  \item Debe generar un campo magnético variable capaz de inducir corriente en el circuito B.
  \item Topología recomendada: oscilador de Colpitts, Hartley o driver MOSFET con feedback resonante.
\end{itemize}

\subsubsection{Circuito receptor B}

\begin{itemize}
  \item Configuración RLC (serie o paralelo), con un LED como carga indicadora.
  \item Sin conexión eléctrica directa con el circuito A.
  \item Debe encender el LED únicamente mediante el acoplamiento inductivo mutuo.
  \item Incorporar condensador variable para ajuste fino de sintonía.
\end{itemize}

% ============================================================
\section{Marco Matemático del Sistema}

\subsection{Frecuencia de resonancia}

La condición de resonancia, que maximiza la impedancia en configuración paralelo y
la minimiza en configuración serie, es:

\begin{formulabox}
\begin{equation}
  f_0 = \frac{1}{2\pi\sqrt{LC}}
  \label{eq:resonancia}
\end{equation}
\end{formulabox}

Para que la transferencia de energía sea máxima, ambos circuitos deben operar en la misma
frecuencia de resonancia $f_0$, de modo que $L_A C_A = L_B C_B$.

\subsection{Factor de calidad y selectividad}

El factor de calidad $Q$ cuantifica la eficiencia energética del circuito resonante y su
selectividad en frecuencia:

\begin{equation}
  Q_{\text{serie}} = \frac{1}{R}\sqrt{\frac{L}{C}}, \qquad
  Q_{\text{paralelo}} = R\sqrt{\frac{C}{L}}
  \label{eq:Qfactor}
\end{equation}

El ancho de banda a $-\SI{3}{\dB}$ del circuito resonante es:
\begin{equation}
  \Delta f_{-3\,\text{dB}} = \frac{f_0}{Q}
  \label{eq:BW}
\end{equation}

Un $Q$ alto proporciona mayor selectividad pero menor tolerancia al desajuste de sintonía.
Se recomienda $Q \in [10, 50]$ para facilitar el ajuste experimental.

\subsection{Acoplamiento mutuo inductivo}

La inductancia mutua~$M$ entre las dos bobinas se relaciona con el coeficiente de
acoplamiento~$k$ mediante:

\begin{formulabox}
\begin{equation}
  M = k\sqrt{L_A L_B}, \qquad k \in [0,\,1]
  \label{eq:mutua}
\end{equation}
\end{formulabox}

Para dos bobinas circulares coaxiales de radios $a_A$, $a_B$ y separadas una distancia~$d$,
la aproximación de Neumann para campo magnético lejano proporciona:

\begin{equation}
  M \approx \frac{\mu_0 \pi N_A N_B\, a_A^2 a_B^2}{2\left(d^2 + a_A^2\right)^{3/2}}
  \qquad (a_B \ll d)
  \label{eq:Neumann}
\end{equation}

\subsection{Eficiencia de transferencia de potencia}

La eficiencia máxima teórica de transferencia de potencia en el punto de operación resonante
está dada por:

\begin{formulabox}
\begin{equation}
  \eta_{\max} = \frac{k^2 Q_A Q_B}{\left(1 + \sqrt{1 + k^2 Q_A Q_B}\right)^2}
  \label{eq:eficiencia}
\end{equation}
\end{formulabox}

De la expresión~\eqref{eq:eficiencia} se desprende que para maximizar $\eta$ deben maximizarse
simultáneamente $k$, $Q_A$ y $Q_B$. La figura de mérito del sistema es el producto $kQ$:
\begin{equation}
  \text{FOM} = k\sqrt{Q_A Q_B}
  \label{eq:FOM}
\end{equation}

\subsection{Inductancia de una bobina helicoidal}

Para una bobina cilíndrica de $N$ espiras, radio $r$ y longitud~$\ell$:
\begin{equation}
  L = \frac{\mu_0 \mu_r N^2 \pi r^2}{\ell} \cdot K_L
  \label{eq:inductancia}
\end{equation}
donde $K_L$ es el factor de corrección de Nagaoka (función de la relación $2r/\ell$), que toma el
valor $K_L \to 1$ para solenoides largos ($\ell \gg 2r$) y decrece para bobinas planas.

\subsection{Tensión inducida en el circuito receptor}

En régimen sinusoidal y a la frecuencia de resonancia:
\begin{equation}
  V_{\text{oc},B} = j\omega_0 M \, I_A = j\omega_0 \, k \sqrt{L_A L_B}\, I_A
  \label{eq:VOC}
\end{equation}

La corriente pico en el transmisor viene determinada por la impedancia de la fuente y la
reactancia del circuito A. En resonancia la impedancia de A es puramente resistiva ($Z_A = R_A$),
luego $I_A = V_s / R_A$.

% ============================================================
\section{Valores de Diseño Orientativos}

La Tabla~\ref{tab:valores} proporciona un punto de partida razonado para el diseño experimental:

\begin{table}[h!]
\centering
\caption{Parámetros de diseño orientativos para el sistema de acoplamiento inductivo}
\label{tab:valores}
\small
\rowcolors{2}{azulclaro!25}{white}
\begin{tabularx}{\linewidth}{@{} l >{\centering\arraybackslash}X >{\centering\arraybackslash}X >{\raggedright\arraybackslash}X @{}}
  \toprule[1.2pt]
  \rowcolor{azulprincipal}
  \color{white}\textbf{Parámetro} &
  \color{white}\textbf{Circuito TX (A)} &
  \color{white}\textbf{Circuito RX (B)} &
  \color{white}\textbf{Observaciones} \\
  \midrule
  Inductancia $L$        & \SI{50}{}--\SI{200}{\micro\henry}  & \SI{50}{}--\SI{200}{\micro\henry}   & Igualar para resonancia común \\
  Capacidad $C$          & \SI{10}{}--\SI{100}{\nano\farad}   & Variable (condensador ajustable)     & $C_B$ ajustable para sintonía fina \\
  Resistencia $R$        & $R_{\text{int}}$ bobina             & $R_{\text{LED}} + R_{\text{serie}}$ & Minimizar $R_A$ para maximizar $Q$ \\
  Frecuencia $f_0$       & \multicolumn{2}{c}{\SI{100}{\kilo\hertz} -- \SI{1}{\mega\hertz}} & Rango cómodo en protoboard \\
  Tensión $V_{CC}$       & \SI{5}{}--\SI{12}{\volt DC}        & ---                                 & Baja tensión de seguridad \\
  N.º de espiras $N$     & 10--50                              & 10--50                              & Ajustar según radio elegido \\
  Radio de bobina $r$    & \SI{3}{}--\SI{8}{\centi\metre}     & \SI{3}{}--\SI{8}{\centi\metre}      & Bobinas planas o cilíndricas \\
  Factor $Q$ estimado    & 20--50                              & 20--50                              & Alambres de cobre, núcleo aire \\
  \bottomrule[1.2pt]
\end{tabularx}
\end{table}

\textbf{Ejemplo de cálculo:} Para $L = \SI{100}{\micro\henry}$ y $f_0 = \SI{500}{\kilo\hertz}$:
\[
  C = \frac{1}{(2\pi f_0)^2 L}
    = \frac{1}{(2\pi \cdot 5\times 10^5)^2 \cdot 10^{-4}}
    \approx \SI{1.01}{\nano\farad}
\]

% ============================================================
\section{Tareas a Realizar --- Parte Práctica}

\subsection{T-01: Diseño de circuitos}

\begin{enumerate}
  \item Calcule los valores óptimos de $L$, $C$ y $R$ para ambos circuitos, estableciendo una
        frecuencia de resonancia común $f_0$ mediante la ecuación~\eqref{eq:resonancia}.
  \item Justifique la elección de configuración RLC (serie o paralelo) en función del
        perfil de la carga (LED) y del objetivo de maximizar la tensión o la corriente inducida.
  \item Determine el factor $Q$ de ambos circuitos mediante~\eqref{eq:Qfactor} y estime
        su influencia sobre la eficiencia y la selectividad.
  \item Calcule la inductancia mutua $M$ para la geometría de bobinas elegida usando~\eqref{eq:Neumann}.
\end{enumerate}

\subsection{T-02: Implementación del sistema}

\begin{enumerate}
  \item Construya ambos circuitos en protoboard o, en su defecto, simúlelos en
        LTspice, Falstad o Multisim. En la simulación deberá reproducir las
        condiciones de acoplamiento mutuo mediante el parámetro~$k$.
  \item Implemente un mecanismo de ajuste de sintonía en el circuito RX:
        \begin{itemize}
          \item Condensador variable (\textit{trimmer}) en paralelo con $C_B$, o
          \item Núcleo de ferrita con posición ajustable dentro de la bobina $L_B$.
        \end{itemize}
  \item Verifique el encendido del LED a diferentes distancias $d$ y orientaciones relativas
        de las bobinas.
\end{enumerate}

\subsection{T-03: Optimización y caracterización experimental}

\begin{enumerate}
  \item Experimente con distintas geometrías de bobina para maximizar $k$:
        bobinas circulares, helicoidales planas (\textit{pancake}) y cilíndricas.
  \item Determine la distancia máxima de operación $d_{\max}$ y represente
        experimentalmente $\eta(d)$.
  \item Cuantifique la eficiencia de transferencia:
        \begin{equation}
          \eta_{\text{exp}} = \frac{P_{\text{out}}}{P_{\text{in}}}
                            = \frac{V_{\text{LED}} \cdot I_{\text{LED}}}{V_{CC} \cdot I_{CC}}
          \label{eq:etaexp}
        \end{equation}
  \item Estudie la degradación de $\eta$ en función del ángulo de desalineamiento~$\alpha$
        entre los ejes de las bobinas y ajuste el modelo teórico a los datos medidos.
\end{enumerate}

% ============================================================
\section{Criterios de Evaluación --- Parte Práctica}

\begin{criteriobox}[title={Rúbrica de evaluación práctica}, fonttitle=\small\bfseries\color{verdemed}]
\begin{enumerate}
  \item \textbf{Funcionalidad del sistema:} encendido reproducible y estable del LED en las condiciones
        de referencia. Se valorará la distancia máxima alcanzada.
  \item \textbf{Eficiencia energética:} eficiencia medida $\eta_{\text{exp}}$ según~\eqref{eq:etaexp},
        maximizada mediante optimización de $k$ y $Q$.
  \item \textbf{Calidad del análisis:} contraste cuantitativo entre predicciones teóricas y medidas
        experimentales. Discusión argumentada de las discrepancias observadas.
  \item \textbf{Documentación técnica:} memoria clara con esquemas de circuito, tablas de medidas,
        gráficas de $\eta(d)$ y conclusiones.
\end{enumerate}
\end{criteriobox}

% ============================================================
\newpage
\section*{Condiciones de Entrega}
\addcontentsline{toc}{section}{Condiciones de Entrega}

\noindent\begin{tabularx}{\linewidth}{@{} >{\bfseries}p{3cm} X @{}}
  \toprule[1.2pt]
  \rowcolor{azulprincipal}
  \color{white}\textbf{Entregable} & \color{white}\textbf{Descripción} \\
  \midrule
  Parte teórica &
    Documento técnico de \textbf{máximo 15 páginas} que incluya: diagrama de bloques,
    desarrollo matemático completo, justificación del diseño de antena, análisis de
    consideraciones prácticas y referencias bibliográficas en formato IEEE. \\[4pt]
  Parte práctica &
    \textbf{Demostración funcional} (en directo o en vídeo) del sistema montado, acompañada
    de una memoria descriptiva con esquemas de circuito, procedimiento de montaje, tabla de
    medidas, gráficas de eficiencia y conclusiones del proceso de optimización. \\[4pt]
  \bottomrule[1.2pt]
\end{tabularx}

\vspace{12pt}
\begin{tcolorbox}[
  enhanced, colback=azulprincipal!8!white, colframe=azulprincipal,
  boxrule=0.5pt, arc=3pt, left=10pt, right=10pt
]
\small
\textbf{Alcance de la evaluación.}
El desarrollo de esta prueba tiene como fin la introducción al contexto de proyecto, a la vez que mostrar \textbf{la capacidad del candidato para aprender y profundizar} 
Se valora especialmente la capacidad para documentar el proceso con rigor técnico y de relacionar los fenómenos físicos, siendo consciente de las limitaciones de los modelos ideales y recurriendo al enfoque modal cuántico cuando sea necesario. 
\end{tcolorbox}


% ============================================================
%  REFERENCIAS
% ============================================================
\newpage
\section*{Referencias Bibliográficas de Referencia}
\addcontentsline{toc}{section}{Referencias Bibliográficas}

\begin{enumerate}[label={[\arabic*]}]
  \item C.~A.~Balanis, \textit{Antenna Theory: Analysis and Design}, 4.ª ed.
        Hoboken, NJ: Wiley, 2016.
  \item D.~M.~Pozar, \textit{Microwave Engineering}, 4.ª ed.
        Hoboken, NJ: Wiley, 2012.
  \item S.~Ramo, J.~R.~Whinnery y T.~Van~Duzer, \textit{Fields and Waves in
        Communication Electronics}, 3.ª ed. New York: Wiley, 1994.
  \item A.~Kurs, A.~Karalis, R.~Moffatt, J.~D.~Joannopoulos, P.~Fisher y
        M.~Solja\v{c}i\'c, ``Wireless power transfer via strongly coupled magnetic
        resonances,'' \textit{Science}, vol.~317, n.º~5834, pp.~83--86, jul.~2007.
  \item W.~X.~Zhong y S.~Y.~R.~Hui, ``Maximum energy efficiency tracking for
        wireless power transfer systems,'' \textit{IEEE Trans.\ Power Electron.},
        vol.~30, n.º~7, pp.~4025--4034, 2015.
  \item H.~Haus y J.~Melcher, \textit{Electromagnetic Fields and Energy}.
        MIT OpenCourseWare, 1989. [En línea]. Disponible en:
        \url{https://ocw.mit.edu}
  \item K.~Finkenzeller, \textit{RFID Handbook: Fundamentals and Applications
        in Contactless Smart Cards and Identification}, 2.ª ed.
        Chichester: Wiley, 2003.
\end{enumerate}

% ─── PIE FINAL ───────────────────────────────────────────
\vfill
\noindent\color{grisbordes}\rule{\linewidth}{0.4pt}
\begin{center}
  \small\color{gray}
  \textit{Prueba Técnica de Ingeniería --- Sistemas de Acoplamiento Electromagnético}\\
  Documento de carácter académico. Todos los derechos reservados.
\end{center}

\end{document}